% sections/analysis.tex

\subsection{问题背景与任务拆解}
微网由光伏、负荷、储能和外购电通道组成，需在预测误差、合同计费和储能约束下平衡常规购电与高价应急购电。相关研究通常将各设备纳入统一的能量管理优化框架\cite{thirunavukkarasu2022optimization,nemati2018unitcommitment}。全年数据结构与季节性特征如图\ref{fig:data_overview}所示。
\begin{figure}[htbp]
\centering
\includegraphics[width=0.85\textwidth]{figures/data-overview.png}
\caption{全年负荷、光伏与电价数据概览及光伏季节性。}
\label{fig:data_overview}
\end{figure}

\subsection{数据画像与预处理}\label{subsec:data-preprocess}
官方全年数据以10分钟为间隔，未发现缺失值或越界记录，故不作插值和删改。功率按$E=P\Delta t$转换为时段电量，其中$\Delta t=1/6\ \mathrm{h}$；第$t$行对应区间$[(t-1)\times10,t\times10)$ min。所有预测只使用决策前已经获得的信息，1月用于预热，2--12月用于正式统计。

基于这些数据，四个问题依次处理“确定性基准—日前计划—日内合同修正—波动电价扩展”四个层次。建模与求解流程如图\ref{fig:route}所示。
\begin{figure}[htbp]
\centering
\includegraphics[width=0.9\textwidth]{figures/route.png}
\caption{全文建模与求解技术流程}
\label{fig:route}
\end{figure}

具体递进关系如下：
\begin{enumerate}
    \item 问题一建立典型日确定性调度基准；
    \item 问题二扩展到仅在每日0:00签约的全年日前调度；
    \item 问题三利用0:00、6:00、12:00和18:00预报调整未执行合同；
    \item 问题四加入波动电价，比较固定合同与可调合同。
\end{enumerate}

各全年方案均从1月1日$6000\ \mathrm{kWh}$的初始储电量启动。1月用于积累历史样本并使储能状态进入连续运行过程，2月1日至12月31日的334天用于正式统计；预热期仍执行完整预测、优化和结算流程。预测数据不足时按“只取已发生观测”降级：$d\ge7$时，F1负荷取$L_{d-7,t}$、光伏取前7日同一时段均值；$1\le d\le6$时取此前已有日的同一时段均值；$d=0$时采用附件1典型日数据。早期残差不足时风险修正项取零，因而全过程均满足非前视要求。

总购电费用是主要评价指标，同时报告应急购电量、闲置合同电量、弃光量和储能终态；预测误差用MAE、RMSE等指标衡量。模型的关键在于：只能使用决策时刻可得的信息，同时处理预测误差与储能跨时段约束。
